\documentclass[11pt]{article}

% Packages for XeLaTeX
\usepackage{fontspec}
\usepackage{amsmath,amssymb,amsthm}
\usepackage{geometry}
\usepackage{hyperref}
\usepackage{enumitem}
\usepackage{mathtools}
\usepackage{bm}
\usepackage{framed}
\usepackage{xcolor}

\geometry{margin=1in}

% Custom commands
\newcommand{\E}{\mathbb{E}}
\newcommand{\R}{\mathbb{R}}
\newcommand{\N}{\mathcal{N}}
\newcommand{\ELBO}{\text{ELBO}}
\newcommand{\tr}{\text{tr}}

\title{ENEL445 Project Planning Report:\\
Optimisation Methods for Variational Inference}
\author{David Ewing (82171165)\\
ENEL445 -- Applied Engineering Optimisation\\
University of Canterbury}
\date{March 17, 2026}

\begin{document}

\maketitle

\section{Background and Problem Intuition}

\subsection{Why Variational Inference?}

Bayesian inference provides principled uncertainty quantification but requires computing intractable posterior distributions. For most real-world models, the posterior integral $p(\bm{\theta} \mid \bm{y}) = p(\bm{y} \mid \bm{\theta})p(\bm{\theta})/p(\bm{y})$ cannot be evaluated analytically. Variational Inference (VI) converts this inference problem into an \textbf{optimisation problem}, making it directly amenable to the gradient-based methods taught in ENEL445.

\subsection{The Core Idea}

Instead of computing the true posterior, VI finds the ``best'' approximation $q(\bm{\theta})$ from a tractable family by maximising the Evidence Lower Bound (ELBO). This transforms intractable integration into tractable optimisation---a natural fit for an optimisation course project.

\section{Optimisation Problem Formulation}

\subsection{Decision Variables}

For Bayesian linear regression with $p$ predictors, the variational parameters (decision variables) are:
\begin{itemize}
\item $\bm{\mu}_\beta \in \R^p$: Posterior mean of regression coefficients
\item $\bm{\Sigma}_\beta \in \R^{p \times p}$: Posterior covariance (symmetric positive definite)
\item $a_e, b_e > 0$: Shape and rate parameters of Gamma distribution for precision
\end{itemize}

\textbf{Total dimension}: $d = p + \frac{p(p+1)}{2} + 2 = 11$ variables for $p=3$ predictors.

\subsection{Objective Function}

Maximize the Evidence Lower Bound (ELBO):
\begin{equation}
\ELBO(\bm{\phi}) = \E_q[\log p(\bm{y}, \bm{\beta}, \tau_e)] - \E_q[\log q(\bm{\beta}, \tau_e)]
\end{equation}

where $\bm{\phi} = (\bm{\mu}_\beta, \bm{\Sigma}_\beta, a_e, b_e)$ are the variational parameters.

\subsection{Constraints}

\begin{itemize}
\item $\bm{\Sigma}_\beta \succ 0$: Covariance matrix must be positive definite
\item $a_e, b_e > 0$: Gamma parameters must be positive
\item Optional: Minimum variance constraints $\sigma^2_{\beta_j} \geq \sigma^2_{\text{min}}$ to prevent underdispersion
\end{itemize}

These constraints can be handled via parameterization (e.g., Cholesky decomposition for $\bm{\Sigma}_\beta$) or penalty methods.

\subsection{Computational Characteristics}

\begin{itemize}
\item Gradient computation: $O(np^2)$ per evaluation (dominated by matrix operations)
\item Hessian computation: $O(np^2 + d^3)$ for Newton's method
\item Parameters exhibit coupling (changing $\bm{\mu}_\beta$ affects optimal $\bm{\Sigma}_\beta$)
\item Must handle special functions: digamma $\psi(\cdot)$, trigamma $\psi'(\cdot)$ for Gamma distribution
\end{itemize}

\section{Test Case: Bayesian Linear Regression}

The test case uses Bayesian linear regression with:
\begin{align}
\bm{y} &\sim \N(\bm{X}\bm{\beta}, \tau_e^{-1} \bm{I}_n) \quad \text{(likelihood)}\\
\bm{\beta} &\sim \N(\bm{0}, \lambda_\beta^{-1} \bm{I}_p) \quad \text{(prior on coefficients)}\\
\tau_e &\sim \text{Gamma}(\alpha_e, \gamma_e) \quad \text{(prior on precision)}
\end{align}

Mean-field approximation:
\begin{align}
q(\bm{\beta}, \tau_e) &= q(\bm{\beta}) \, q(\tau_e)\\
q(\bm{\beta}) &= \N(\bm{\mu}_\beta, \bm{\Sigma}_\beta)\\
q(\tau_e) &= \text{Gamma}(a_e, b_e)
\end{align}

\textbf{Variational parameters}: $\bm{\phi} = (\bm{\mu}_\beta, \bm{\Sigma}_\beta, a_e, b_e)$ with dimension $d = p + \frac{p(p+1)}{2} + 2$

For $p=3$ predictors: $d = 3 + 6 + 2 = 11$ parameters to optimise.

\section{Overview of Optimisation Methods}

\subsection{Baseline Method: CAVI}

\textbf{Coordinate Ascent Variational Inference (CAVI)} is VI-specific and not a standard ENEL445 method. It updates each variational parameter block sequentially using closed-form solutions:

\begin{align}
\bm{\Sigma}_\beta &\leftarrow \left(\frac{a_e}{b_e} \bm{X}^T\bm{X}\right)^{-1}\\
\bm{\mu}_\beta &\leftarrow \frac{a_e}{b_e} \bm{\Sigma}_\beta \bm{X}^T\bm{y}\\
a_e &\leftarrow \alpha_e + \frac{n}{2}\\
b_e &\leftarrow \gamma_e + \frac{1}{2}\left[\|\bm{y} - \bm{X}\bm{\mu}_\beta\|^2 + \tr(\bm{X}^T\bm{X}\bm{\Sigma}_\beta)\right]
\end{align}

\textbf{Properties}: Monotonic ELBO increase, linear convergence, $O(np^2)$ per iteration. Provides analytical benchmark for validation.

\subsection{Course-Based Methods}

The following methods from ENEL445 apply directly to maximizing ELBO

\subsection{Mapping Course Methods to VI Problem}

\begin{enumerate}[leftmargin=*]
\item \textbf{Gradient Ascent} (Course: Pages 27-28, gradient descent)
\begin{itemize}
\item Standard iterative update: $\bm{\phi}^{(t+1)} = \bm{\phi}^{(t)} + \alpha_t \nabla \ELBO(\bm{\phi}^{(t)})$
\item Line search from Chapter 3.2 (bisection/Wolfe conditions)
\item Stopping criteria: $\|\nabla \ELBO\| < \epsilon$
\item Convergence: Linear rate
\end{itemize}

\item \textbf{Newton's Method} (Course: Equation 3.26)
\begin{itemize}
\item Standard Newton update: $\bm{\phi}^{(t+1)} = \bm{\phi}^{(t)} + \alpha_t [\nabla^2 \ELBO]^{-1} \nabla \ELBO$
\item Requires computing and inverting the Hessian matrix
\item Computational cost: $O(d^3)$ per iteration for Hessian inversion
\item Convergence: Quadratic rate near optimum
\item Challenge: Hessian may not be positive definite away from optimum
\end{itemize}

\item \textbf{BFGS Quasi-Newton} (Course: Equations 4.50-4.51, one of five required methods from Project 2)
\begin{itemize}
\item BFGS update formula builds inverse Hessian approximation
\item Uses gradient differences: $\bm{s}_k = \bm{\phi}^{(k)} - \bm{\phi}^{(k-1)}$, $\bm{y}_k = \nabla \ELBO(\bm{\phi}^{(k)}) - \nabla \ELBO(\bm{\phi}^{(k-1)})$
\item Woodbury identity for efficient updates of inverse Hessian
\item Computational cost: $O(d^2)$ per iteration
\item Convergence: Superlinear rate
\item In practice, this is the workhorse method: avoids storing full Hessian but captures essential curvature information
\end{itemize}
\end{enumerate}

\subsection{Why These Methods?}

\begin{itemize}
\item \textbf{Gradient Ascent}: Baseline first-order method, expected to be slowest but most robust
\item \textbf{Newton's Method}: Best convergence rate (quadratic) if Hessian well-conditioned
\item \textbf{BFGS}: ``Best general-purpose optimiser'' for smooth problems (course recommendation)
\item \textbf{CAVI}: Domain-specific baseline exploiting problem structure
\end{itemize}

Comparison will reveal whether general optimisation methods can compete with or outperform problem-specific algorithms

\section{Solution Plan}

\subsection{Phase 1: Data Generation and Baseline (Completed)}

\textbf{Data simulation}:
\begin{itemize}
\item Generate $n = 100$ observations, $p = 3$ predictors
\item Design matrix $\bm{X} \sim \N(0, 1)$, true coefficients $\bm{\beta}_{\text{true}} = [1, -0.5, 2]^T$
\item Noise: $\epsilon \sim \N(0, 0.5^2)$
\end{itemize}

\textbf{Mathematical foundations (completed)}:
\begin{itemize}
\item Complete ELBO derivation for Bayesian linear regression
\item CAVI closed-form update equations
\item Gradient expressions: $\nabla_{\bm{\mu}_\beta} \ELBO$, $\nabla_{\bm{\Sigma}_\beta} \ELBO$, $\nabla_{a_e} \ELBO$, $\nabla_{b_e} \ELBO$
\item Hessian block structure: diagonal and mixed blocks
\item BFGS update formulation with Woodbury identity
\end{itemize}

All derivations documented in \texttt{vi\_optimization\_solution.tex} (13 pages)

\subsection{Phase 2: Implementation (In Progress)}

\begin{enumerate}[leftmargin=*]
\item \textbf{Data Generation} (5\%)
\begin{itemize}
\item Simulate data: $n = 100$ observations, $p = 3$ predictors
\item Generate $\bm{X}$ from standard normal, $\bm{\beta}_{\text{true}} = [1, -0.5, 2]^T$
\item Add noise: $\epsilon \sim \N(0, 0.5^2)$
\end{itemize}

\item \textbf{Baseline Implementation: CAVI} (15\%)
\begin{itemize}
\item Implement closed-form coordinate updates
\item Track ELBO convergence
\item Stopping criterion: $|\ELBO^{(t+1)} - \ELBO^{(t)}| < 10^{-6}$
\end{itemize}

\item \textbf{Method 1: Gradient Ascent} (25\%)
\begin{itemize}
\item Implement ELBO gradient computation
\item Line search: bisection method or Wolfe conditions
\item Stopping criterion: $\|\nabla \ELBO\| < 10^{-6}$
\item Track iteration count and ELBO values
\end{itemize}

\item \textbf{Method 2: Newton's Method} (25\%)
\begin{itemize}
\item Implement Hessian computation (block structure)
\item Regularization: add $\lambda \bm{I}$ if Hessian not positive definite
\item Line search for step size $\alpha_t$
\item Track convergence and computational cost
\end{itemize}

\item \textbf{Method 3: BFGS Quasi-Newton} (25\%)
\begin{itemize}
\item Initialise inverse Hessian approximation: $\bm{B}_0 = \bm{I}_d$
\item Implement BFGS update formula with Woodbury identity
\item Line search with Wolfe conditions
\item Track superlinear convergence
\end{itemize}

\item \textbf{Comparison and Analysis} (5\%)
\begin{itemize}
\item Compare convergence speed (iterations to $10^{-6}$ tolerance)
\item Compare computational cost (time per iteration, total time)
\item Compare solution quality (posterior approximation accuracy)
\item Analyse when each method performs best
\end{itemize}
\end{enumerate}

\subsection{Success Metrics}

\begin{enumerate}
\item \textbf{Convergence}: All methods reach ELBO optimum within tolerance
\item \textbf{Correctness}: Solutions match CAVI baseline (analytical benchmark)
\item \textbf{Performance hierarchy}: 
\begin{itemize}
\item Expected: BFGS $>$ Newton $>$ Gradient Ascent (in convergence speed)
\item CAVI likely fastest for this conjugate problem
\end{itemize}
\item \textbf{Understanding}: Clear explanation of why each method performs as it does
\end{enumerate}

\section{Progress To Date}

\subsection{Completed Work}

\begin{enumerate}
\item \textbf{Problem formulation}: Full optimisation problem specified with decision variables, objective function (ELBO), and constraints

\item \textbf{Mathematical derivations}: 13-page technical document containing:
\begin{itemize}
\item Complete ELBO derivation from first principles
\item CAVI algorithm with closed-form updates
\item Analytical gradient expressions for all 11 variational parameters
\item Hessian block computations for Newton's method
\item BFGS quasi-Newton formulation
\item Penalty method formulation for constrained optimisation
\end{itemize}

\item \textbf{Method selection justification}: Demonstrated direct mapping between course material (Chapters 3-4) and VI optimisation:
\begin{itemize}
\item Gradient ascent $\leftrightarrow$ Algorithm 3.2 (gradient descent)
\item Newton's method $\leftrightarrow$ Equation 3.26
\item BFGS $\leftrightarrow$ Equations 4.50-4.51 (one of five required project methods)
\end{itemize}

\item \textbf{Validation strategy}: Established CAVI as analytical benchmark for correctness checking
\end{enumerate}

\subsection{Preliminary Results}

While full implementation is pending, the mathematical analysis reveals:

\begin{itemize}
\item \textbf{Gradient structure}: $\nabla_{\bm{\mu}_\beta} \ELBO \propto \bm{X}^T(\bm{y} - \bm{X}\bm{\mu}_\beta)$ resembles least-squares gradient, suggesting ELBO surface is well-behaved

\item \textbf{Hessian properties}: Diagonal blocks are negative definite (ELBO is locally concave), supporting Newton's method applicability

\item \textbf{Dimensionality}: With $d=11$ parameters, Hessian inversion costs $\approx 1,300$ flops---manageable for Newton's method without L-BFGS reduction

\item \textbf{Parameter coupling}: Mixed Hessian blocks non-zero, indicating gradient-based methods exploiting joint updates should outperform coordinate ascent
\end{itemize}

\subsection{Next Steps}

\begin{enumerate}
\item Implement CAVI baseline in Python (Week 7)
\item Implement gradient ascent with line search (Week 7)
\item Implement Newton's method with regularization (Week 8)
\item Implement BFGS quasi-Newton (Week 8)
\item Monte Carlo comparison: convergence iterations, wall-clock time, ELBO values (Week 8-9)
\item Final report with performance analysis (Week 12)
\end{enumerate}

\section{Alignment with Course Requirements}

This problem hits all the major learning objectives from ENEL445. I need to formulate an unconstrained optimisation problem (done), compute gradients analytically and numerically (done), and implement three methods from Class with line search and convergence criteria. The comparison directly tests course concepts: when does Newton beat BFGS? When is gradient descent actually fast enough? These answers depend on the problem structure, which is the whole point.

\textbf{Key insight}: This isn't a toy problem. VI is how people actually do Bayesian inference in production, so the methods that win here are the ones that matter.

\section{Timeline}

\begin{tabular}{ll}
\textbf{March 17, 2026} & Mathematical derivations completed \\
\textbf{March 18, 2026} & Data generation and CAVI implementation \\
\textbf{March 19, 2026} & Gradient ascent and Newton's method implementation \\
\textbf{March 20, 2026} & BFGS implementation and comparison analysis \\
\textbf{March 20, 2026} & \textbf{Progress report submission} \\
\end{tabular}

\section{Summary}

I've formulated a Bayesian inference problem as an unconstrained optimisation task with 11 decision variables and a well-defined ELBO objective. The mathematics is done. What remains is to code four methods—CAVI baseline, gradient ascent, Newton, BFGS—and show which converges fastest and why. The comparison should be instructive: CAVI exploits problem structure but doesn't scale; Newton is fast locally but needs a good starting point; BFGS sits in the middle, practical and robust.

\end{document}
